\FloatBarrier
\subsection{Attribute Program Synthesis}
\label{sec:synthesis}
The mapping procedure establishes which logical axes occur in source tensors
and where they occur. Synthesis uses these relationships to fill the
replacement's attribute holes. It normalizes value expressions, identifies
which axes each replacement operation must preserve or reduce, and queries
mappings to construct its attributes.

Figure~\ref{fig:auto-matching} previews this process for distribution of a dot
over addition. The source is lowered into Axes IR and normalized; the RHS
arithmetic skeleton supplies reads to match before its attributes are known.
Read correspondence allows reuse of source coordinates, after which axis
positions construct the two missing dot descriptors. Here addition has no
broadcasting and preserves paired axis order, so both inferred descriptors
equal the source descriptor. The figure's reads $A[I,K]$, $B[I,K]$, and
$C[K,J]$ abbreviate group coordinates, rather than asserting that groups are
contiguous in physical layouts. The skeleton's question marks denote unresolved
reduction information. Equivalence is checked only after construction.

\paperfigure{1}{figures/Auto attribute matching_ from rules to attributes.pdf}{Matching normalized operand reads supplies coordinate correspondences for constructing replacement attributes.}{fig:auto-matching}{Three stages connect a distribution rule and operation semantics to normalized
terms, read matching, position queries, and completed replacement descriptors.}

For attribute reconstruction beyond descriptor reuse, continue the guarded
reassociation example under an arithmetic policy permitting the required laws:
\begin{equation}
 \underbrace{\operatorname{dot}(\operatorname{dot}(X,Y),Z)}_{\mathrm{LHS}}
 \ \longrightarrow\
 \underbrace{\operatorname{dot}(X,\operatorname{dot}(Y,Z))}_{\mathrm{RHS}}.
\end{equation}
The rule supplies the replacement graph. Let $T=\operatorname{dot}(X,Y)$ be
the source intermediate, $U=\operatorname{dot}(Y,Z)$ the replacement
intermediate, and $O,O'$ their respective final results.

\FloatBarrier
\subsubsection{Value Expression Normalization}
\label{sec:normalization}
The graphs combine their reads in different orders. Normalization identifies
corresponding computations before constructing the replacement's attributes.
With output coordinates fixed and indexed reads abbreviated, the LHS is
\begin{equation}
 e_{\mathrm{LHS}}=\sum_L\left(\sum_K XY\right)Z.
\end{equation}
Each sum ranges over the coordinates of its group. The abbreviated notation
hides, but does not discard, the mappings established in
Section~\ref{sec:axis-mapping}.

TEPL uses a sum-of-products representation similar to SPORES's R-polynomial
form~\cite[Section~2.3, Definition~2.4]{spores}. A term contains a coefficient,
indexed operand reads, and symbolic reduction axes:
\begin{equation}
 t=(c,[r_1,\ldots,r_m],R),\qquad
 \operatorname{value}(t)=\sum_R c\prod_{i=1}^{m}r_i.
 \label{eq:normalized-term}
\end{equation}
Each read retains its tensor identity and coordinate expression. $R$ is an
ordered sequence of symbolic axes, and an empty sequence denotes no summation.
The empty product is $1$, permitting constants. A normalized expression is a
list $N=[t_1,\ldots,t_n]$ denoting
\begin{equation}
 \operatorname{value}(N)=\sum_{j=1}^{n}\operatorname{value}(t_j).
\end{equation}
Reads within a term multiply, separate terms add, and the empty term list
represents zero. Examples, again abbreviating indexed reads, are:
\begin{center}
\begin{tabular}{@{}ll@{}}\toprule
Expression & Normalized term list\\\midrule
$XY$ & \code{[(1, [X, Y], [])]}\\
$X(Y+Z)=XY+XZ$ & \code{[(1, [X, Y], []), (1, [X, Z], [])]}\\
$2XY-Z$ & \code{[(2, [X, Y], []), (-1, [Z], [])]}\\
$\sum_K XY$ & \code{[(1, [X, Y], K)]}\\\bottomrule
\end{tabular}
\end{center}
This similarity to SPORES concerns representation; it does not transfer a
completeness result to TEPL's supported reasoning.

Under the permitted multiplication and reduction laws, the example normalizes
to
\begin{equation}
 \sum_L\left(\sum_K XY\right)Z
 =\sum_L\sum_K XYZ
 =\sum_{L,K}XYZ.
\end{equation}
The inner reduction contributes reads \code{X}, \code{Y}, and axes $K$.
Multiplication incorporates \code{Z} after checking that it does not depend
on $K$. The outer reduction prepends $L$ after establishing that $L$ and $K$
are disjoint. Thus
\begin{codeblock}[H]
\begin{lstlisting}
Value expression:
Reduce(L,
  Product(
    Reduce(K, Product(X, Y)),
    Z))

Normalized term:
t = (1, [X, Y, Z], L ++ K)
\end{lstlisting}
\end{codeblock}

Unknown RHS attributes mean its actual reduction axes and read coordinates
are not yet available. Synthesis instead normalizes its \emph{arithmetic
skeleton}: it expands declared value operations while temporarily omitting
axis information. The skeleton for $X(YZ)$ has coefficient $1$, reads
\code{[X,Y,Z]}, and an empty reduction sequence. That empty sequence is a
property of the skeleton representation, not a claim that the completed RHS
has no reductions.

Synthesis matches each RHS read with the same operand's LHS read and reuses
its coordinates. For several terms, it first matches terms with the same
operands. Multiple possible read-coordinate assignments must be shown to
agree; otherwise synthesis remains unresolved. This correspondence provides
coordinates for construction. Reduction domains, coefficients, and full
value equivalence still have to be checked after attributes are known.

\FloatBarrier
\subsubsection{Determine Which Axes Each Operation Preserves}
\label{sec:keep-axes}
After read matching, the logical axes in each RHS input are known, but their
roles in the replacement may change. The source inner dot contracts $K$;
the source outer dot contracts $L$, partitioned into $\Isel$ and $\Jsel$.
The final output retains $B,\Irem,\Jrem,F$. Source mappings show that
\code{Y} contains $B,K,\Jsel,\Jrem$ and \code{Z} contains
$B,\Isel,\Jsel,F$. These lists specify membership rather than physical order.
To build $U$, the compiler must classify these axes as batch, contraction,
left-free, or right-free. In particular, $K$ must survive $U$ for a later
contraction with \code{X}.

\begin{figuretext}{figures/Infer Axis Roles in U.pdf}{An intermediate preserves axes needed outside its computation or in the final output and reduces the remaining shared axes.}{fig:infer-u}{A flow compares axes inside the replacement intermediate with axes needed
outside it and in the output, retaining required axes and contracting the rest.}
Figure~\ref{fig:infer-u} shows the decision for an unbatched instance.
Starting with axes in \code{Y} or \code{Z}, it first preserves those needed
by \code{X}, then those required by the final output. The remaining shared
group, $\Jsel$, can be reduced inside $U$. The figure omits batch group $B$;
the equations below retain it for the guarded general case.

\paragraph{General rule for preserving axes.}
An intermediate must keep an axis used to compute it if that axis is also
needed outside the computation or in the final output. For intermediate $t$,
its replacement subtree includes the operation and its input reads. In $U$,
reads of \code{Y} and \code{Z} are inside; the read of \code{X} is outside.
Let $\Axes(x)$ denote the set of logical axes in tensor $x$'s interface,
ignoring order. In a matched product, let $\Axes_{\mathrm{in}}(t)$ and
$\Axes_{\mathrm{out}}(t)$ collect axes used by reads inside and outside the
subtree. With original output $O$, the required set is
\begin{equation}
\begin{aligned}
 \Keep(t)&=\Axes_{\mathrm{in}}(t)\cap\\
 &\quad\bigl(\Axes_{\mathrm{out}}(t)\\
 &\qquad{}\cup\Axes(O)\bigr).
\end{aligned}
 \label{eq:keep-axes}
\end{equation}
These are sets of logical coordinates, with order determined later. For
multiple matched products, the required sets must agree.
Reducing an axis removes its coordinate from the intermediate. An outside
read using that coordinate must still pair its values with the intermediate
at matching coordinates; an output axis must remain available to distinguish
output elements. These requirements give the union in
Equation~\eqref{eq:keep-axes}.
\end{figuretext}

Intersecting with the inside set excludes axes
used only outside $t$, which need not pass through it. Axes used only inside
can be reduced there if the operation's declared semantics permits it.

For $U$, the inside set is $\Axes(Y)\cup\Axes(Z)$ and the outside set is
$\Axes(X)$. Using group names as shorthand for all their constituent axes,
\begin{equation}
\begin{aligned}
 \Axes_{\mathrm{in}}(U)
    &=\{B,K,\Isel,\Jsel,\Jrem,F\},\\
 \Axes_{\mathrm{out}}(U)
    &=\{B,K,\Isel,\Irem\},\\
 \Axes(O)&=\{B,\Irem,\Jrem,F\},\\
 \Keep(U)&=\{B,K,\Isel,\Jrem,F\}.
\end{aligned}
\end{equation}
Group $K$ occurs inside and outside and so is preserved. $\Jsel$ is needed
only inside and can be contracted. Although $\Irem$ belongs to the final
output, it occurs outside $U$ and reaches the final dot through \code{X}.
The resulting decisions are summarized in Table~\ref{tab:keep-decisions}.
At the final dot, $K$ and $\Isel$ are shared by \code{X} and $U$ but absent
from $O$, so they are reduced together.

\begin{table}[htbp]
\caption{Preservation decisions for groups present in $Y$ or $Z$.}
\label{tab:keep-decisions}
\begin{tabularx}{\linewidth}{@{}lYY@{}}\toprule
Group & Need outside $U$ or in $O$ & Action in $U$\\\midrule
$B$ & Used by $X$ and present in $O$ & Preserve as batch axes.\\
$K$ & Used by $X$ & Preserve for final contraction.\\
$\Isel$ & Used by $X$; paired axes occur in $Z$ & Preserve for final contraction.\\
$\Jsel$ & Absent from $X$ and $O$ & Contract between $Y$ and $Z$.\\
$\Jrem,F$ & Present in $O$ & Preserve as free axes.\\\bottomrule
\end{tabularx}
\end{table}

\paragraph{From required axes to operation groups.}
The required set identifies surviving coordinates. Assigning coordinates to
operation groups additionally uses which operands contain them and whether
the declared value expression reduces the group. For an operation producing
$t$ from $x_1,\ldots,x_m$, let $\Axes_i=\Axes(x_i)$, and let $E_g$ be the
nonempty set of operands bound to declared group $g$. Axes appearing in exactly
those operands form
\begin{equation}
 D_g=\left(\bigcap_{i\in E_g}\Axes_i\right)
     \setminus\left(\bigcup_{i\notin E_g}\Axes_i\right).
\end{equation}
The candidate assignment is
\begin{equation}
 \Axes_g=\begin{cases}
 D_g\setminus\Keep(t),&\text{if the value definition reduces }g,\\
 D_g\cap\Keep(t),&\text{otherwise.}
 \end{cases}
 \label{eq:group-assignment}
\end{equation}
This equation assigns membership, with ordered positions constructed afterward.
For a dot with operands $a,b$, shared axes in $\Keep(t)$ are batch, and shared
axes outside $\Keep(t)$ are contracted. Left-only and right-only axes in
$\Keep(t)$ form the respective free groups. Thus the assignments are:
\begin{center}
\small
\begin{tabular}{@{}lcccc@{}}\toprule
Operation & Batch & Contracted & Left free & Right free\\\midrule
$U=\operatorname{dot}(Y,Z)$ & $B$ & $\Jsel$ & $K,\Jrem$ & $\Isel,F$\\
$O'=\operatorname{dot}(X,U)$ & $B$ & $K,\Isel$ & $\Irem$ & $\Jrem,F$\\\bottomrule
\end{tabular}
\end{center}
Free axes follow operand order. The rule relies on distinguishable declared
group roles; ambiguous assignments remain unresolved. Groups introduced by an
explicit output domain rather than operand bindings are handled separately
in Section~\ref{sec:synthesis-algorithm}.

\FloatBarrier
\subsubsection{Construct Attributes from the Mapping Table}
\label{sec:attribute-construction}
Group membership says which axes a dot uses, but its attributes need ordered
positions in actual operands. A query $\pos(P,G)$ returns $p$ satisfying
$\gath(P,p)=G$, using existing mappings and their compositions. Both sides
of a paired binding query the same \emph{ordered} group, ensuring that
corresponding entries select paired axes.

\begin{figuretext}{figures/Constructing Dot Attributes from Axis Mappings.pdf}{Paired position queries construct dot attributes; the completed intermediate supplies mappings for the final dot.}{fig:attribute-construction}{Paired queries on Y and Z construct the inner dot, an interface bridge
establishes U's layout, and paired queries on X and U construct the final dot.}
Figure~\ref{fig:attribute-construction} shows the two stages for the unbatched
example. The first pair of queries constructs the inner replacement's
contraction lists. Its completed interface $P_U$ then supplies mappings for
the final dot's queries. In the figure, $G$ abbreviates
$\Isel\catseq K$, and preserved-group chips specify membership, not order.
The short fields \code{lhs_contract} and \code{rhs_contract} abbreviate the
contraction attributes. The following derivation also includes batch lists.


\paragraph{1. Construct the inner replacement dot.}
For $U=\operatorname{dot}(Y,Z)$, the source table already locates $B$ and
$\Jsel$ in both operands. Its attributes are
\begin{equation}
\begin{aligned}
 b_L^U&=b_Y,&b_R^U&=b_Z,\\
 c_L^U&=\gath(j_Y,v_J),&c_R^U&=z_J.
\end{aligned}
\end{equation}
By Equation~\eqref{eq:z-partners}, $z_J=\gath(c_R^{\mathrm{outer}},u_J)$,
so both lists select $\Jsel$ in the same order. Free positions are
\begin{equation}
\begin{aligned}
 f_Y^U&=\\
 &\ocomp(\rho(P_Y),b_Y\catseq c_L^U),\\
 f_Z^U&=\\
 &\ocomp(\rho(P_Z),b_Z\catseq c_R^U),
\end{aligned}
\end{equation}
and the dot's output order constructs
\begin{equation}
\begin{aligned}
 P_U&=B\catseq\gath(P_Y,f_Y^U)\\
 &\catseq\gath(P_Z,f_Z^U).
\end{aligned}
 \label{eq:u-interface}
\end{equation}
The compiler extends mappings to $P_U$ using this layout. This is necessary
because $K$ and $\Jrem$ can be interleaved in \code{Y}, as can $\Isel$ and
$F$ in \code{Z}. The interface preserves those physical orders instead of
assuming that each surviving group is contiguous.

\end{figuretext}

\paragraph{2. Construct the final replacement dot.}
The final dot contracts $K$ and $\Isel$. Filtering source reduction order
$L\catseq K$ to those axes gives $\Isel\catseq K$: axes from $I$ occur
within $L$, before $K$. Use that same order for both contraction lists:
\begin{equation}
\begin{aligned}
 c_L^{O'}&=\gath(i_X,v_I)\catseq k_X,\\
 c_R^{O'}&=\pos(P_U,\Isel)\catseq\pos(P_U,K).
\end{aligned}
\end{equation}
A fixed binding determines group order, and inferred paired lists must follow
it. Otherwise the implementation prefers source reduction order for reduced
groups and source output order for retained groups, appending remaining axes
in operand order. Table~\ref{tab:replacement-attributes} completes the two
descriptors.
\begin{table}[htbp]
\caption{Candidate attributes for the replacement dots.}
\label{tab:replacement-attributes}
\begin{tabularx}{\linewidth}{@{}lYY@{}}\toprule
Attribute & $U=\operatorname{dot}(Y,Z)$ & $O'=\operatorname{dot}(X,U)$\\\midrule
Left batch & $b_Y$ & $b_X$\\
Right batch & $b_Z$ & $\rho(B)$\\
Left contraction & $\gath(j_Y,v_J)$ & $\gath(i_X,v_I)\catseq k_X$\\
Right contraction & $z_J$ & $\pos(P_U,\Isel)$\newline $\catseq\pos(P_U,K)$\\\bottomrule
\end{tabularx}
\end{table}
Removing the final contraction axes leaves
\begin{equation}
 P_{O'}=B\catseq\Irem\catseq\Jrem\catseq F=P_O.
\end{equation}
These attributes therefore have the candidate pairings and output order.
Operation legality and value equivalence still require checking.

\FloatBarrier
\subsubsection{Produce and Verify Attribute Programs}
\label{sec:emit-programs}
The tables contain symbolic group queries, while a rule application must
compute concrete attributes. The compiler resolves each query to expressions
over source attributes and shapes. For the inner replacement's right
contraction list, for example,
\begin{equation}
\begin{aligned}
 u_J&=\mpo\bigl(c_L^{\mathrm{outer}},
                   |B|+|I|+\rho(J)\bigr)_1,\\
 c_R^U&=\gath(c_R^{\mathrm{outer}},u_J),
\end{aligned}
\label{eq:source-only-example}
\end{equation}
where subscript $1$ selects the first returned sequence. Group lengths and
ranges are themselves derived from source metadata. This expression constructs
positions at each match without accessing tensor elements. Queries involving
$P_U$ similarly substitute its layout and compose mappings until only source
metadata remains. Emitted programs contain neither unresolved RHS attributes
nor logical-group queries.

For calculated accesses, position mappings alone need not determine every
attribute. The implementation additionally matches supported index expressions,
including coordinate coefficients and offsets. Composed slices, for example,
have the candidate attributes
\begin{equation}
\begin{aligned}
 \mathrm{start}'&=\mathrm{start}_{\mathrm{inner}}
   +\mathrm{stride}_{\mathrm{inner}}\,\mathrm{start}_{\mathrm{outer}},\\
 \mathrm{stride}'&=\mathrm{stride}_{\mathrm{inner}}\,
                  \mathrm{stride}_{\mathrm{outer}}.
\end{aligned}
\end{equation}
These expressions must also satisfy resulting access bounds. Coupled or
inverse equations outside the supported fragment remain unresolved.

\paragraph{Check the completed replacement.}
With attributes filled, the RHS has actual coordinates and reduction domains.
In the example, $U$ reduces $\Jsel$ and the final dot reduces
$\Isel\catseq K$:
\begin{equation}
\begin{aligned}
 e_{\mathrm{RHS}}
 &=\sum_{\Isel,K}X\left(\sum_{\Jsel}YZ\right)\\
 &=\sum_{\Isel,K,\Jsel}XYZ.
\end{aligned}
\end{equation}
The selected subgroups partition $L$, and the mappings preserve coordinate
pairing. The normalized RHS can therefore be compared to
$\sum_{L,K}XYZ$ under the permitted reduction laws. Verification also checks
output order, shape, dtype, and operation legality.

After construction, a fresh verification context reconstructs the completed
rule using the emitted programs. Required conditions must follow from legal
source metadata and the declared guard; construction obligations do not
become new assumptions. On success, the result is a collection of checked
source-dependent programs, compiled once and evaluated on each matched
graph's metadata. Unestablished construction or verification relationships
leave the rule unresolved.

\FloatBarrier
\subsubsection{General Synthesis Algorithm}
\label{sec:synthesis-algorithm}
The procedure uses each operation's declared groups, value expression, and
output order, rather than requiring the replacement to contain dots. Its
inputs are a rule's source graph, replacement graph with attribute holes,
operation definitions, and declared conditions. It returns checked attribute
programs or an unresolved result.

The state retains $(V,M,\Pi)$ and adds a table $C$ of proposed attribute
expressions. \todo{Confirm operation-qualified keys for $C$: the pseudocode
uses an attribute name as shorthand, but different replacement operations
can have fields with the same name.} In the pseudocode, \code{.axes} and \code{.value} denote the two
components of $V[t]$. Read correspondence supplies required output sets,
and replacement nodes are processed after their inputs. The construction
invariant requires every completed node to have resolved legal attributes,
an interface and value expression consistent with them, and mappings
satisfying Equation~\eqref{eq:mapping-invariant}. Whole-rule equivalence is
checked after completion.

\begin{paperalgorithm}[H]
\caption{Synthesize attribute programs.}\label{alg:synthesis}
\begin{lstlisting}[language=Python]
def synthesize_attribute_programs(rule):
    V, M, partitions = build_symbolic_axis_mapping(rule.lhs)
    source = normalize(V[rule.lhs].value)
    skeleton = normalize(arithmetic_skeleton(rule.rhs))
    match = correspond(source, skeleton)
    required = required_output_axes(match, V[rule.lhs].axes)
    bind_replacement_reads(V, match)
    C = {}

    for op in producer_first(rule.rhs):
        if op.is_input:
            continue
        inputs = [V[t] for t in op.inputs]
        if op.has_attribute_holes:
            groups = select_groups(op.definition, inputs, required[op])
            groups = order_groups(groups, op, source, inputs)
            for binding in op.unresolved_axis_bindings:
                P = inputs[binding.operand].axes
                G = groups[binding.group]
                C[binding.attribute] = positions(P, G, M, partitions)
            complete_other_attributes(op, inputs, groups, match, C)

        completed = substitute_attributes(op, C)
        V[op] = lower_and_check(completed, inputs, M, partitions)

    check_equivalence(V[rule.lhs], V[rule.rhs])
    programs = source_only_expressions(C)
    verify_in_fresh_context(rule, programs)
    return programs
\end{lstlisting}
\end{paperalgorithm}

\code{correspond} matches normalized reads, and \code{required_output_axes}
computes $\Keep(t)$ for each replacement intermediate. \code{select_groups}
uses operand membership and reduction roles as in
Section~\ref{sec:keep-axes}. Explicit output-domain groups are separate: they
come from required output axes absent from every operand, and their sizes
determine domain attributes. Ambiguous group assignments remain unresolved.
\code{order_groups} respects fixed bindings and otherwise follows source
and operand orders. \code{positions} abbreviates mapping lookup and position
construction for these ordered groups.

\code{complete_other_attributes} handles supported index coefficients,
offsets, output domains, and output permutations. \code{lower_and_check}
checks each completed operation and publishes its interface and mappings for
consumers. \code{check_equivalence} compares fully lowered values, reduction
domains, output order, shape, and dtype. \code{source_only_expressions}
resolves constructors to source metadata, and the final check reconstructs
the completed rule in a fresh context. The pseudocode omits failure paths:
unsupported or unproved steps return an unresolved result. Successful output
is a program for every missing attribute, ready for match-time evaluation.
